\appendix
\section*{Appendix: Reproducibility and Evidence Ledger}

The paper package is structured so that a reviewer can trace every main-text
number to a copied source table or figure. The Gate 5 CSV and the successor
branch-summary CSV (comma-separated-value) files in
\texttt{results\_tables/} provide the numerical values
reported in Tables \ref{tab:gate5-results} and \ref{tab:successor-branch}.
The copied figure registry records the title, evidence status, source data,
generator, digest, and claim boundary for every figure included here. The
figures retain their registry identifiers in their captions.

The release package provides scripts, fixed configuration records, data cards,
and decision logs. Reproduction of the published development
figures should begin from the tagged repository state and should preserve the
lock on calibration and final-test data. Re-running a report-generation script
is different from authorizing new fitting or a new experimental campaign.

\subsection*{Interpretation Checklist for Reviewers}

The following checklist summarizes the intended reading of the paper:
\begin{enumerate}
\item Treat Gate 3 as an eligible-scope simulator consistency result, not a
mission validation result.
\item Treat Gate 5 as rejection of the Q01 fidelity-kernel hypothesis on this
development benchmark, not a universal statement that QML cannot work.
\item Treat the projected Q01b kernel, feasibility-only FQK classifier, and 16
successors D034--D049 as development-only negative evidence that cannot revise
Gate 5 or open Gate 6.
\item Treat D046 as evidence that an apparent improvement over the C06
physics-residual model must also be tested against a like-for-like classical
control.
\item Treat the recall-first analysis as a future lesson about retaining
feasible plans while controlling infeasible admissions, not as a deployed
safety filter.
\item Do not infer quantum advantage, hardware performance, propellant saving,
or flight readiness from any figure or table in this package.
\end{enumerate}
